\subsection*{Grounded answer writing across typefaces}
The final two glyphs of every formatted question in the 1,788-question
reading assay were identical. The most frequent complete reference bitmap
sequence occurred six times. A writer observing only that suffix can match
at most six reference answers. Removing the common cue and grouping the
raw questions by their last two visible signatures raises the empirical
ceiling to 184. Neither calculation involves a new fitted model.

We retained the trained passage reader and supplied its selection to a writer
trained to draw Serif strokes from Sans images. On 3,170 official development
questions, the reader selected a correct span 565 times. With fine local
source detail, the writer reproduced 124 complete Serif answers exactly;
with the same reading intervals and only a coarse layout it reproduced one
(Fig.~\ref{fig:grounded-detail}). Direct copies of the selected Sans images
matched zero Serif answers. The difference between detailed and coarse writing
was 3.88 percentage points (95\% passage interval 3.23 to 4.57).
Exchanging questions within passages gave nine exact answers to 3,158
original questions, a paired decline of 3.61 points (2.94 to 4.31).
Seven of 3,100 cyclic pairs requiring distinct answers had both complete
outputs correct. These are conditional transcriptions of selected evidence.
The source of answer length is the learned span boundary.

The writer reproduced 2,600 of 4,742 familiar nonempty training pairs exactly.
On 537 withheld shapes it reproduced none, with mean ink F1 .660 against .724
for direct copying. A 10,833-parameter convolutional fallback raised F1 to
.751, an improvement of .0271 over copying (95\% shape interval .0231 to
.0312), below the prespecified .03 point requirement. On 314 further Han
shapes absent from every training passage the fallback reached .7565 against
.7334 for copying. Neither set yielded an exact new glyph. Repainting every
development answer retained the same 124 complete correct responses and
seven both-correct pairs.
